% ----------------------------------------------------------------------
%  1.8. Кристаллография мартенситных превращений
% ----------------------------------------------------------------------
\begin{orientir}
Мартенситное превращение~— перестройка решётки \vazhno{сдвигом}, без
диффузии. В~1.8.1 по рельефу на шлифе выясним, что макроскопически это
\vazhno{деформация с инвариантной плоскостью} и почему, кроме деформации
решётки, нужна ещё деформация, не меняющая решётку (скольжение или
двойникование). В~1.8.2~— как ГЦК-аустенит переходит в
мартенсит (деформация Бейна) и какие \vazhno{ориентационные соотношения}
наблюдаются на опыте.
\end{orientir}

Название «мартенсит» (по имени немецкого металловеда А.~Мартенса)
первоначально относилось к твёрдой составляющей закалённой стали. Сейчас
мартенситные (их называют также \emph{сдвиговыми}) превращения известны в
сплавах, чистых металлах, керамиках, минералах, неорганических
соединениях, затвердевших газах и полимерах.

\begin{opredelenie}
\term{Мартенситное превращение}~— бездиффузионное фазовое превращение, при
котором решётка перестраивается \emph{кооперативным} смещением атомов на
расстояния меньше межатомного; каждый атом сохраняет своих соседей, а
состав фаз одинаков.
\end{opredelenie}

Превращение идёт с очень большой скоростью: кристалл мартенсита в стали
растёт со скоростью порядка $10^5$~см/с ($\sim1$~км/с), сравнимой со
скоростью звука. Как правило, мартенсит образуется \emph{атермически}:
количество мартенсита определяется температурой охлаждения ниже точки
$M_\text{н}$ и почти не растёт при выдержке (хотя известно и изотермическое
мартенситное превращение).

\begin{masshtab}[мартенсит]{-10/-9.6/{смещения атомов}, -8.3/-7.3/{внутренние двойники}, -7/-4/{рейки и пластины}}
Каждый атом смещается меньше чем на межатомное расстояние (на десятые
доли нанометра и меньше), но согласованно с миллиардами соседей.
Кристаллы мартенсита~— рейки толщиной 0,1--0,5~мкм в низкоуглеродистой
стали и пластины размером вплоть до зерна (10--100~мкм) в
высокоуглеродистой; внутренние двойники в пластинах~— 5--50~нм.
Пластина длиной 100~мкм при скорости роста $\sim1$~км/с вырастает за
$\sim10^{-7}$~с.
\end{masshtab}

\subsection{Морфология мартенситных превращений}\label{ssec:martensite-morph}

Важнейшие сведения о геометрии превращения дал металлографический
анализ \emph{рельефа}, возникающего на предварительно полированной
поверхности шлифа там, где образовался кристалл мартенсита
(рис.~\ref{fig:mart-relief}). Установлено:
\begin{itemize}
\item поверхность превращённой области наклоняется к плоскости шлифа,
но остаётся \emph{плоской};
\item прямые риски, нанесённые до превращения, остаются \emph{непрерывными}
и \emph{прямыми}, но испытывают излом там, где выходит на поверхность
граница кристалла мартенсита~— \term{плоскость габитуса}.
\end{itemize}

\begin{figure}[ht]
\centering
\begin{tikzpicture}[cubeproj, scale=1.35]
  \def\dz{0.55}
  % исходная фаза (x от 0 до 2)
  \fill[fmBlueBg] (0,0,0) -- (2,0,0) -- (2,0,1.5) -- (0,0,1.5) -- cycle;
  \fill[fmBlueBg!70!white] (0,0,1.5) -- (2,0,1.5) -- (2,2.5,1.5) -- (0,2.5,1.5) -- cycle;
  % мартенсит (x от 2 до 4): верхняя грань наклонена
  \fill[fmAmberBg] (2,0,0) -- (4,0,0) -- (4,0,1.5+\dz) -- (2,0,1.5) -- cycle;
  \fill[fmAmber!25] (2,0,1.5) -- (4,0,1.5+\dz) -- (4,2.5,1.5+\dz) -- (2,2.5,1.5) -- cycle;
  \fill[fmAmber!15] (4,0,0) -- (4,2.5,0) -- (4,2.5,1.5+\dz) -- (4,0,1.5+\dz) -- cycle;
  \draw[atomLine] (0,0,0) -- (4,0,0) -- (4,0,1.5+\dz) -- (2,0,1.5) -- (0,0,1.5) -- cycle
    (0,0,1.5) -- (0,2.5,1.5) -- (2,2.5,1.5) -- (4,2.5,1.5+\dz) -- (4,0,1.5+\dz)
    (4,0,0) -- (4,2.5,0) -- (4,2.5,1.5+\dz);
  \draw[atomLine!50, dashed] (0,2.5,1.5) -- (0,2.5,0) -- (4,2.5,0) (0,2.5,0) -- (0,0,0);
  % плоскость габитуса ABC
  \fill[fmRed!20, opacity=0.6] (2,0,0) -- (2,0,1.5) -- (2,2.5,1.5) -- (2,2.5,0) -- cycle;
  \draw[fmRed!80!black, line width=0.9pt] (2,0,0) -- (2,0,1.5) -- (2,2.5,1.5);
  \draw[fmRed!80!black, dashed] (2,2.5,1.5) -- (2,2.5,0) -- (2,0,0);
  \node[font=\small, fmRed!80!black, anchor=north] at (2,0,0) {$A$};
  \node[font=\small, fmRed!80!black, anchor=east] at (2,0,1.5) {$B$};
  \node[font=\small, fmRed!80!black, anchor=south] at (2,2.5,1.5) {$C$};
  % риски
  \foreach \y in {0.6,1.25,1.9}{
    \draw[fmBlue, line width=1pt] (0,\y,1.5) -- (2,\y,1.5) -- (4,\y,1.5+\dz);}
  \node[font=\small] at (1,0,0.6) {1};
  \node[font=\small] at (3,0,0.6) {2};
  \node[font=\small, fmBlue] at (0.6,1.9,1.75) {3};
\end{tikzpicture}
\caption{Рельеф на поверхности шлифа при образовании кристалла
мартенсита: 1~— исходная фаза; 2~— мартенсит; 3~— прямые риски,
нанесённые на поверхность до превращения; $ABC$~— плоскость габитуса.
Риски остаются прямыми и непрерывными, но изламываются на следе
плоскости габитуса}
\label{fig:mart-relief}
\end{figure}

Это означает, что рельеф создан \emph{однородной деформацией с
инвариантной плоскостью}: существует плоскость (габитус), которая не
искажается и не поворачивается, а все смещения направлены в одну
сторону и пропорциональны расстоянию от неё. Если $\vb p$~— единичная
нормаль к инвариантной плоскости, а $\vb d$~— вектор смещения точек,
отстоящих от неё на единицу длины, то
\begin{keyeq}\label{eq:ips}
\vb r'=\vb r+\vb d\,(\vb p\cdot\vb r),\qquad P=I+\vb d\,\tilde{\vb p}
\end{keyeq}
(ср. со смещением при скольжении~\eqref{eq:slip-displacement}). Два
предельных случая (рис.~\ref{fig:ips}): $\vb d\parallel\vb p$~—
одноосное растяжение или сжатие (дилатация), $\vb d\perp\vb p$~— простой
сдвиг. Любую деформацию с инвариантной плоскостью можно представить как
их комбинацию. У мартенсита в сталях сдвиговая составляющая $\approx0{,}2$, а
нормальная~— несколько процентов (удельный объём растёт на 3--4\,\%).

\begin{figure}[ht]
\centering
\begin{tikzpicture}[scale=1.4]
  \foreach \X/\h/\s/\l in {0/1.35/0/а, 2.6/1.0/0.5/б, 5.2/1.3/0.4/в}{
    \begin{scope}[shift={(\X,0)}]
      \fill[fmAmberBg] (0,0) -- (1.4,0) -- (1.4+\s,\h) -- (\s,\h) -- cycle;
      \draw[atomLine, dashed] (0,0) rectangle (1.4,1);
      \draw[fmAmber!80!black, line width=0.8pt] (0,0) -- (1.4,0) -- (1.4+\s,\h) -- (\s,\h) -- cycle;
      \draw[fmRed, line width=2pt] (-0.25,0) -- (1.65,0);
      \draw[arr, fmBlue] (0.7,1) -- ({0.7+\s},\h);
      \node[font=\small] at (0.7,-0.35) {\textit{\l}};
    \end{scope}}
  \node[font=\footnotesize, fmBlue, anchor=west] at (0.75,1.18) {$\vb d$};
  \node[font=\footnotesize, fmRed, anchor=west] at (7.05,0) {инвариантная плоскость};
\end{tikzpicture}
\caption{Деформации с инвариантной плоскостью (красная): \textit{а}~—
дилатация ($\vb d\parallel\vb p$); \textit{б}~— сдвиг ($\vb d\perp\vb p$);
\textit{в}~— сдвиг с дилатацией. Пунктир~— исходная форма}
\label{fig:ips}
\end{figure}

Если перестройку решётки представлять как деформацию с инвариантной
плоскостью, то граница между когерентными областями новой и исходной фаз
и есть инвариантная плоскость~— плоскость габитуса. Чтобы сохранить
контакт с непревращённой частью кристалла, деформация должна включать не
только чистую деформацию (растяжение, сжатие, сдвиг), но и \emph{поворот}.

\paragraph{Деформация решётки и деформация с инвариантной решёткой}
Опыт показал, что макроскопическая деформация превращённого объёма
(найденная по рельефу) \emph{отличается} от деформации, переводящей
решётку аустенита в решётку мартенсита. Разницу обеспечивает
дополнительная деформация, которая не меняет решётку,~— \term{деформация
с инвариантной решёткой}: скольжение или двойникование
(рис.~\ref{fig:lis}). Деформация решётки задаёт \emph{структуру} новой фазы,
а деформация с инвариантной решёткой~— её \emph{форму}.

\begin{figure}[ht]
\centering
\begin{tikzpicture}[scale=0.95]
  \def\w{0.42} \def\h{0.38} \def\g{0.105}
  % (а) исходная решётка
  \begin{scope}
    \foreach \j in {0,...,7}{\foreach \i in {0,...,3}{
      \draw[fmBlue, fill=fmBlueBg] ({\i*\w},{\j*\h}) rectangle ++(\w,\h);}}
    \node[font=\small] at (0.84,-0.4) {\textit{а}};
  \end{scope}
  % (б) деформация решётки: каждая ячейка скошена, форма изменилась
  \begin{scope}[shift={(2.6,0)}]
    \draw[atomLine!50, dashed] (0,0) rectangle ({4*\w},{8*\h});
    \foreach \j in {0,...,7}{\foreach \i in {0,...,3}{
      \draw[fmAmber!80!black, fill=fmAmberBg] ({\i*\w+\j*\g},{\j*\h}) -- ++(\w,0) -- ++(\g,\h) -- ++(-\w,0) -- cycle;}}
    \node[font=\small] at (1.2,-0.4) {\textit{б}};
  \end{scope}
  % (в) + скольжение: каждые 4 слоя сдвинуты назад на период
  \begin{scope}[shift={(5.6,0)}]
    \draw[atomLine!50, dashed] (0,0) rectangle ({4*\w},{8*\h});
    \foreach \j in {0,...,7}{
      \pgfmathsetmacro{\off}{\j*\g-\w*floor(\j/4)}
      \foreach \i in {0,...,3}{
        \draw[fmAmber!80!black, fill=fmAmberBg] ({\i*\w+\off},{\j*\h}) -- ++(\w,0) -- ++(\g,\h) -- ++(-\w,0) -- cycle;}}
    \draw[fmRed, line width=1pt] (-0.15,{4*\h}) -- ({4*\w+0.6},{4*\h});
    \node[font=\small] at (1.0,-0.4) {\textit{в}};
  \end{scope}
  % (г) + двойникование: чередование слоёв с противоположным скосом
  \begin{scope}[shift={(8.6,0)}]
    \draw[atomLine!50, dashed] (0,0) rectangle ({4*\w},{8*\h});
    \foreach \j in {0,...,7}{
      \pgfmathtruncatemacro{\m}{mod(\j,4)}
      \ifnum\m<2
        \pgfmathsetmacro{\off}{\m*\g}
        \foreach \i in {0,...,3}{
          \draw[fmAmber!80!black, fill=fmAmberBg] ({\i*\w+\off},{\j*\h}) -- ++(\w,0) -- ++(\g,\h) -- ++(-\w,0) -- cycle;}
      \else
        \pgfmathsetmacro{\off}{(4-\m)*\g}
        \foreach \i in {0,...,3}{
          \draw[fmRed!70!black, fill=fmRedBg] ({\i*\w+\off},{\j*\h}) -- ++(\w,0) -- ++(-\g,\h) -- ++(-\w,0) -- cycle;}
      \fi}
    \node[font=\small] at (0.95,-0.4) {\textit{г}};
  \end{scope}
\end{tikzpicture}
\caption{Деформация решётки и деформация с инвариантной решёткой:
\textit{а}~— исходная решётка; \textit{б}~— деформация решётки меняет и
структуру, и форму (пунктир~— исходная форма); \textit{в}~— та же
структура, но скольжение (красная линия~— плоскость скольжения)
возвращает макроскопическую форму, оставляя ступеньки; \textit{г}~— то же
достигается двойникованием: слои двух двойниковых ориентаций (оранжевые и
розовые) чередуются, и средняя форма не меняется}
\label{fig:lis}
\end{figure}

Так объясняется \emph{субструктура} мартенсита: в низкоуглеродистых сталях
это пакеты реек с высокой плотностью дислокаций (деформация с
инвариантной решёткой~— скольжением), в высокоуглеродистых и в сплавах
Fe--Ni~— линзовидные пластины с тонкими внутренними двойниками
(двойникованием).

\subsection{Кристаллография мартенситного превращения в сталях}\label{ssec:martensite-cryst}

При закалке стали ГЦК-решётка аустенита (твёрдого раствора углерода в
$\gamma$-Fe) превращается в \emph{объёмноцентрированную тетрагональную}
(ОЦТ) решётку мартенсита. Как это возможно при минимальных смещениях
атомов, первым показал Бейн.

\paragraph{Соответствие решёток по Бейну}
В двух соседних ГЦК-ячейках можно выделить ОЦТ-ячейку
(рис.~\ref{fig:bain},\,\textit{а}): её основание~— квадрат со стороной
$a_\gamma/\sqrt2$, построенный на направлениях $[110]_\gamma$ и
$[\br110]_\gamma$, высота~— $a_\gamma$ вдоль $[001]_\gamma$, а в центре~—
атом из центра общей грани. Отношение $c/a$ этой ячейки $\sqrt2=1{,}41$.
Чтобы получить ОЦК (или слабо тетрагональную) ячейку мартенсита, её
нужно сжать вдоль $c$ и растянуть в основании~— это и есть
\term{деформация Бейна}:
\begin{keyeq}\label{eq:bain}
B=\begin{pmatrix}\eta_1&0&0\\0&\eta_1&0\\0&0&\eta_3\end{pmatrix}
\approx\begin{pmatrix}1{,}12&0&0\\0&1{,}12&0\\0&0&0{,}83\end{pmatrix},\qquad
\eta_1=\frac{\sqrt2\,a_\alpha}{a_\gamma},\quad \eta_3=\frac{c_\alpha}{a_\gamma}
\end{keyeq}
(в осях $[100]_\alpha$, $[010]_\alpha$, $[001]_\alpha$): растяжение примерно
на 12\,\% вдоль $[100]_\alpha$ и $[010]_\alpha$ и сжатие примерно на 17\,\%
вдоль $[001]_\alpha$. Объём при этом увеличивается в
$\eta_1^2\eta_3\approx1{,}04$ раза.

\begin{figure}[ht]
\centering
\begin{tikzpicture}[cubeproj, scale=2.0]
  % (а) две ГЦК-ячейки и выделенная ОЦТ-ячейка
  \draw[atomLine!45, dashed] (0,1,0) -- (2,1,0) (0,1,0) -- (0,0,0) (0,1,0) -- (0,1,1) (1,1,0) -- (1,1,1);
  \draw[atomLine] (0,0,0) -- (2,0,0) -- (2,1,0) -- (2,1,1) -- (0,1,1) -- (0,0,1) -- cycle
    (0,0,1) -- (2,0,1) -- (2,1,1) (2,0,0) -- (2,0,1) (1,0,0) -- (1,0,1) -- (1,1,1);
  \foreach \p in {(1,0,0),(1.5,0.5,0),(1,1,0),(0.5,0.5,0)}{\draw[fmRed!80!black, line width=1.1pt] \p -- ++(0,0,1);}
  \draw[fmRed!80!black, line width=1.1pt] (1,0,0) -- (1.5,0.5,0) -- (1,1,0) -- (0.5,0.5,0) -- cycle
    (1,0,1) -- (1.5,0.5,1) -- (1,1,1) -- (0.5,0.5,1) -- cycle;
  \foreach \x in {0,1,2}{\foreach \y in {0,1}{\foreach \z in {0,1}{
    \node[atom=2.6mm] at (\x,\y,\z) {};}}}
  \foreach \x in {0.5,1.5}{\foreach \z in {0,1}{\node[atom=2.6mm] at (\x,0.5,\z) {};}
    \foreach \y in {0,1}{\node[atom=2.6mm] at (\x,\y,0.5) {};}}
  \foreach \x in {0,1,2}{\node[atom=2.6mm] at (\x,0.5,0.5) {};}
  \foreach \p in {(1,0,0),(1.5,0.5,0),(1,1,0),(0.5,0.5,0),(1,0,1),(1.5,0.5,1),(1,1,1),(0.5,0.5,1)}{
    \node[atomB=2.6mm] at \p {};}
  \node[atomI=2.6mm] at (1,0.5,0.5) {};
  \node[font=\small] at (1,0,-0.35) {\textit{а}};
  % (б) деформация Бейна
  \begin{scope}[shift={(2.9,0,0)}]
    \def\a{0.707}
    \draw[fmRed!80!black, line width=0.8pt] (0,0,0) -- (\a,0,0) -- (\a,\a,0) (0,0,0) -- (0,0,1) -- (\a,0,1) -- (\a,\a,1) -- (0,\a,1) -- (0,0,1) (\a,0,0) -- (\a,0,1) (\a,\a,0) -- (\a,\a,1);
    \draw[fmRed!50, dashed] (0,0,0) -- (0,\a,0) -- (\a,\a,0) (0,\a,0) -- (0,\a,1);
    \foreach \p in {(0,0,0),(\a,0,0),(\a,\a,0),(0,0,1),(\a,0,1),(\a,\a,1),(0,\a,1),(0,\a,0)}{\node[atomB=2.4mm] at \p {};}
    \node[atomI=2.4mm] at ({\a/2},{\a/2},0.5) {};
    \node[font=\scriptsize, anchor=north] at ({\a/2},0,0) {$a_\gamma/\sqrt2$};
    \node[font=\scriptsize, anchor=east] at (0,0,0.5) {$a_\gamma$};
    \draw[arrB] (1.05,0.35,0.5) -- (1.55,0.35,0.5);
    \node[font=\scriptsize, above] at (1.3,0.35,0.55) {$B$};
    \begin{scope}[shift={(1.85,0,0.08)}]
      \def\b{0.79} \def\c{0.83}
      \draw[atomLine, line width=0.8pt] (0,0,0) -- (\b,0,0) -- (\b,\b,0) (0,0,0) -- (0,0,\c) -- (\b,0,\c) -- (\b,\b,\c) -- (0,\b,\c) -- (0,0,\c) (\b,0,0) -- (\b,0,\c) (\b,\b,0) -- (\b,\b,\c);
      \draw[atomLine!50, dashed] (0,0,0) -- (0,\b,0) -- (\b,\b,0) (0,\b,0) -- (0,\b,\c);
      \foreach \p in {(0,0,0),(\b,0,0),(\b,\b,0),(0,0,\c),(\b,0,\c),(\b,\b,\c),(0,\b,\c),(0,\b,0)}{\node[atomB=2.4mm] at \p {};}
      \node[atomI=2.4mm] at ({\b/2},{\b/2},{\c/2}) {};
      \node[font=\scriptsize, anchor=north] at ({\b/2},0,0) {$a_\alpha$};
      \node[font=\scriptsize, anchor=west] at (\b,0,{\c/2}) {$c_\alpha$};
      \draw[arr, fmBlue] ({\b/2},{\b/2},{\c+0.45}) -- ({\b/2},{\b/2},{\c+0.12});
      \draw[arr, fmBlue] (-0.45,0,{\c/2}) -- (-0.12,0,{\c/2});
    \end{scope}
    \node[font=\small] at (1.5,0,-0.35) {\textit{б}};
  \end{scope}
\end{tikzpicture}
\caption{Соответствие решёток по Бейну: \textit{а}~— две ГЦК-ячейки
аустенита и выделенная в них ОЦТ-ячейка ($c/a=\sqrt2$; оранжевые атомы~—
её вершины, красный~— центр); \textit{б}~— деформация Бейна: сжатие вдоль
$c$ примерно на 17\,\% и растяжение в основании на 12\,\% переводят её в
ячейку мартенсита}
\label{fig:bain}
\end{figure}

Соответствие Бейна сразу объясняет \emph{тетрагональность} мартенсита.
Атомы углерода в аустените занимают октаэдрические пустоты~— середины
рёбер и центр ГЦК-ячейки (п.~\ref{ssec:elem-types}). При соответствии
Бейна все эти пустоты попадают в \emph{одну} из трёх подрешёток
октаэдрических пустот ОЦТ-решётки~— ту, у которой сплюснутый октаэдр
вытянут вдоль оси $c$. Внедрённые атомы углерода раздвигают решётку вдоль
$c$, и мартенсит оказывается тетрагональным:
$c/a\approx1+0{,}046\,C$, где $C$~— содержание углерода в массовых
процентах.

Деформации Бейна соответствуют \term{ориентационные соотношения Бейна}:
\[
(001)_\gamma\parallel(001)_\alpha,\qquad [100]_\gamma\parallel[110]_\alpha .
\]
В чистом виде они наблюдаются лишь при превращении в очень тонких плёнках.

\paragraph{Деформация Бейна не имеет инвариантной плоскости}
Применим деформацию Бейна к сфере единичного радиуса: она превратится в
эллипсоид с полуосями $\eta_1,\eta_1,\eta_3$ (рис.~\ref{fig:bain-cone}).
Неискажёнными (сохранившими длину) остаются лишь направления, по которым
эллипсоид пересекает сферу,~— они образуют \emph{конус} $L'$, в который
перешёл конус $L$ исходных направлений. Найдём его: точка эллипсоида
$(x,z)$ лежит на сфере, если $x^2+z^2=1$ и
$x^2/\eta_1^2+z^2/\eta_3^2=1$. При $\eta_1=1{,}12$, $\eta_3=0{,}83$
угол конуса $L'$ с осью $c$ равен $56{,}2^\circ$, а исходного конуса $L$~—
$47{,}9^\circ$: неискажённые направления \emph{поворачиваются}.

\begin{figure}[ht]
\centering
\begin{tikzpicture}[scale=1.9]
  \fill[fmBlueBg] (0,0) circle (1);
  \fill[fmAmber!25, opacity=0.8] (0,0) ellipse (1.12 and 0.83);
  \draw[atomLine, line width=0.7pt] (0,0) circle (1);
  \draw[fmAmber!80!black, line width=0.9pt] (0,0) ellipse (1.12 and 0.83);
  \draw[arr, atomLine!70] (-1.35,0) -- (1.45,0) node[right, font=\small, black]{$[100]_\alpha$};
  \draw[arr, atomLine!70] (0,-1.15) -- (0,1.3) node[above, font=\small, black]{$[001]_\alpha$};
  \foreach \sx in {1,-1}{\foreach \sz in {1,-1}{
    \draw[fmBlue, dashed, line width=0.8pt] (0,0) -- ({\sx*0.7417*1.2},{\sz*0.6707*1.2});
    \draw[fmRed, line width=1pt] (0,0) -- ({\sx*0.8307*1.2},{\sz*0.5567*1.2});
    \fill[fmBlue] ({\sx*0.7417},{\sz*0.6707}) circle (0.025);
    \fill[fmRed] ({\sx*0.8307},{\sz*0.5567}) circle (0.025);}}
  \draw[arr, fmGray] (0.7417,0.6707) to[bend left=20] (0.8307,0.5567);
  \node[font=\small, fmBlue, anchor=south] at ({0.7417*1.2},{0.6707*1.2}) {$L$};
  \node[font=\small, fmRed, anchor=west] at ({0.8307*1.2},{0.5567*1.2}) {$L'$};
  \node[font=\footnotesize, anchor=north west] at (1.0,-0.05) {$\eta_1$};
  \node[font=\footnotesize, anchor=south west] at (0.02,0.83) {$\eta_3$};
  \node[font=\footnotesize, fmBlue!70!black, anchor=east] at (-0.72,0.75) {сфера};
  \node[font=\footnotesize, fmAmber!80!black, anchor=east] at (-1.1,-0.35) {эллипсоид};
\end{tikzpicture}
\caption{Деформация Бейна в сечении, содержащем ось $c$: сфера
единичного радиуса переходит в эллипсоид. Неискажённые направления
лежат на конусе $L'$ (пересечение сферы и эллипсоида); до деформации они
лежали на конусе $L$. Плоскости, которая осталась бы и неискажённой, и
неповёрнутой, нет}
\label{fig:bain-cone}
\end{figure}

Пусть плоскость $P'$ после деформации проходит через два неискажённых
направления $a'$ и $b'$ конуса $L'$; до деформации она проходила через
направления $a$ и $b$ (плоскость $P$). Но углы между неискажёнными
направлениями не сохраняются, поэтому \emph{ни одна} плоскость не
остаётся неискажённой. Причина в главных удлинениях: инвариантная
плоскость возможна только если одно из них равно 1, а два других~— по
разные стороны от 1. У деформации Бейна два удлинения больше 1, одно
меньше, и ни одно не равно 1.

\paragraph{Как получается инвариантная плоскость}
Добавление к деформации Бейна обычного скольжения не помогает: скольжение
сопровождается поворотом кристалла, и инвариантная плоскость по-прежнему
не возникает. Феноменологическая теория мартенситного превращения
(Векслер, Либерман, Рид; Боулс, Маккензи) строит наблюдаемую
макроскопическую деформацию $P_1$ из трёх частей:
\begin{keyeq}\label{eq:wlr}
P_1=R\,P\,B ,
\end{keyeq}
где $B$~— деформация Бейна (задаёт решётку мартенсита), $P$~—
деформация с инвариантной решёткой (двойникование или скольжение по
подходящей системе), $R$~— поворот. Величину сдвига $P$ подбирают так,
чтобы у произведения $PB$ одно главное удлинение стало равно 1; тогда
появляется неискажённая плоскость, а поворот $R$ возвращает её в исходное
положение~— получается инвариантная плоскость, т.\,е. плоскость
габитуса. Геометрически: одно неискажённое направление выбирают в
плоскости двойникования $K_1$ (оно сохраняется при двойниковании и
переходит в неискажённое при деформации Бейна), второе~— в плоскости
$K_2$; плоскость, проходящая через них, после двойникования и деформации
Бейна поворачивается, и поворот $R$ совмещает её с исходной. Теория
предсказывает плоскость габитуса, ориентационное соотношение и
макроскопический сдвиг; все три хорошо согласуются с опытом.

\paragraph{Наблюдаемые ориентационные соотношения}
Опыт даёт ориентационные соотношения, близкие к соотношению Бейна, но
отличающиеся от него поворотом $R$ на несколько градусов
(табл.~\ref{tab:or}).

\begin{table}[ht]
\centering
\caption{Ориентационные соотношения при мартенситных превращениях}
\label{tab:or}
\small
\begin{tabular}{@{}lllcl@{}}
\toprule
Соотношение & Плоскости & Направления & Вариантов & Где наблюдается\\
\midrule
Бейна & $(001)_\gamma\parallel(001)_\alpha$ & $[100]_\gamma\parallel[110]_\alpha$ & 3 & тонкие плёнки\\
Курдюмова--Закса & $(111)_\gamma\parallel(011)_\alpha$ & $[10\br1]_\gamma\parallel[11\br1]_\alpha$ & 24 & стали, 0,5--1,4\,\% C\\
Нисиямы & $(111)_\gamma\parallel(011)_\alpha$ & $[\br1\br12]_\gamma\parallel[0\br11]_\alpha$ & 12 & сплавы Fe--Ni\\
Бюргерса & $\plfam{110}_\beta\parallel(0001)_\alpha$ & $\nfam{111}_\beta\parallel\nfam{2\br1\br10}_\alpha$ & 12 & Zr, Ti (ОЦК$\to$ГПУ)\\
\bottomrule
\end{tabular}
\end{table}

В массивных образцах сталей с 0,5--1,4\,\% (масс.) углерода выполняются
соотношения \term{Курдюмова--Закса}: плотноупакованная плоскость $\plfam{111}$
аустенита параллельна плотноупакованной плоскости $\plfam{110}$ мартенсита,
а плотноупакованные направления $\nfam{110}_\gamma$ и $\nfam{111}_\alpha$
в этих плоскостях параллельны. Плоскость габитуса лежит в области
$\plfam{225}_\gamma$--$\plfam{259}_\gamma$. В сплавах Fe--Ni выполняются
соотношения \term{Нисиямы} (Нисиямы--Вассермана)\ispr{В оригинале
направления в соотношении Нисиямы записаны как
$[112]_\gamma\parallel[11\br1]_\alpha$, т.\,е. $\nfam{112}_\gamma\parallel\nfam{111}_\alpha$.
Это неверно: в соотношении Нисиямы--Вассермана
$\nfam{112}_\gamma\parallel\nfam{110}_\alpha$ (эквивалентно
$\nfam{110}_\gamma\parallel\nfam{100}_\alpha$). Оно отличается от
соотношения Курдюмова--Закса поворотом на $5{,}26^\circ$ вокруг нормали к
$(111)_\gamma$.}: плоскости те же, но в плоскости $(111)_\gamma$ решётка
повёрнута на $5{,}26^\circ$, так что параллельны $\nfam{112}_\gamma$ и
$\nfam{110}_\alpha$; габитус около $\plfam{3\,10\,15}_\gamma$.

\begin{smysl}
Ориентационное соотношение~— «отпечаток» механизма превращения:
плотноупакованные плоскости и направления двух фаз совмещены, потому что
решётка перестраивается сдвигом с минимальными смещениями атомов. Число
вариантов (24 для соотношения Курдюмова--Закса) показывает, сколько
по-разному ориентированных кристаллов мартенсита может возникнуть в одном
зерне аустенита.
\end{smysl}

Мартенситный характер имеют и некоторые полиморфные превращения. Так, при
переходе ОЦК $\beta$-Zr в ГПУ $\alpha$-Zr (и в титане) выполняются
\term{соотношения Бюргерса}: $\plfam{110}_\beta\parallel(0001)_\alpha$,
$\nfam{111}_\beta\parallel\nfam{2\br1\br10}_\alpha$. Более подробный анализ
требует рассмотрения дислокационных механизмов, кинетики и термодинамики
мартенситных превращений.

\begin{itogi}
\begin{itemize}
\item Мартенситное превращение~— бездиффузионная кооперативная
перестройка решётки сдвигом; скорость роста $\sim1$~км/с.
\item Макроскопически~— деформация с инвариантной плоскостью (габитусом):
$\vb r'=\vb r+\vb d(\vb p\cdot\vb r)$; на поверхности возникает рельеф.
\item Полная деформация $P_1=RPB$: деформация решётки $B$ (Бейн:
$\eta_1\approx1{,}12$, $\eta_3\approx0{,}83$), деформация с инвариантной
решёткой $P$ (скольжение или двойникование) и поворот $R$.
\item Деформация Бейна сама по себе не имеет инвариантной плоскости
(лишь конус неискажённых направлений).
\item Соответствие Бейна объясняет тетрагональность мартенсита:
$c/a\approx1+0{,}046\,C$.
\item Ориентационные соотношения: Курдюмова--Закса
$\plfam{111}_\gamma\parallel\plfam{110}_\alpha$,
$\nfam{110}_\gamma\parallel\nfam{111}_\alpha$; Нисиямы; Бюргерса
($\beta\to\alpha$ в Zr, Ti).
\end{itemize}
\end{itogi}

\begin{proverka}
\item Какие наблюдения рельефа показывают, что мартенситное превращение~—
деформация с инвариантной плоскостью?
\item Чем деформация решётки отличается от деформации с инвариантной
решёткой? Приведите примеры последней.
\item Найдите по~\eqref{eq:bain} изменение объёма при превращении аустенита
($a_\gamma=0{,}3584$~нм) в мартенсит ($a=0{,}2850$, $c=0{,}2958$~нм).
\item Почему деформация Бейна не имеет инвариантной плоскости?
\item Проверьте, что направление $[10\br1]_\gamma$ лежит в плоскости
$(111)_\gamma$, а $[11\br1]_\alpha$~— в плоскости $(011)_\alpha$.
\end{proverka}
